% !TEX root = ../Thesis.tex

\chapter{基于深度学习的多模光纤散斑图像重建方法}\label{chap:method}

{
本章详细介绍本文提出的基于深度学习的多模光纤散斑图像重建方法，这是本论文的核心内容。

\section{引言}

% TODO: 介绍本章的主要内容和创新点


\section{实验系统搭建}

\subsection{光学系统设计}

% TODO: 介绍实验光路设计，附示意图

\subsection{多模光纤选择}

% TODO: 介绍所使用的多模光纤的参数

\subsection{数据采集系统}

% TODO: 介绍数据采集的硬件和软件系统


\section{数据集构建}

\subsection{数据采集方法}

% TODO: 详细描述如何采集训练数据

\subsection{数据集规模与特征}

% TODO: 介绍构建的数据集的规模、特点

\subsection{数据增强策略}

% TODO: 介绍采用的数据增强方法


\section{改进的网络模型}

\subsection{网络架构设计}

% TODO: 介绍改进的网络架构（如果有改进的话）
% TODO: 参考 ref/个人科研成果 中的论文内容

\subsection{网络创新点}

% TODO: 详细说明相比标准UNet的改进之处

\subsection{损失函数设计}

% TODO: 介绍损失函数的设计（如果有特殊设计）


\section{训练与优化}

\subsection{训练细节}

% TODO: 详细介绍训练过程的各项设置

\subsection{训练过程分析}

% TODO: 展示训练曲线，分析训练过程


\section{实验结果}

\subsection{图像重建结果}

% TODO: 展示不同场景下的重建结果（大量图片）

\subsection{定量性能评估}

% TODO: 使用多种指标进行定量评估（表格）

\subsection{与现有方法的对比}

% TODO: 与其他方法进行对比实验


\section{鲁棒性分析}

\subsection{对光纤扰动的鲁棒性}

% TODO: 分析方法对光纤弯曲、扭转等扰动的鲁棒性

\subsection{对噪声的鲁棒性}

% TODO: 分析方法对噪声的鲁棒性

\subsection{泛化能力分析}

% TODO: 分析方法的泛化能力


\section{实时性分析}

% TODO: 分析重建速度，讨论实时应用的可能性


\section{讨论}

% TODO: 讨论实验结果，分析方法的优缺点


\section{本章小结}

% TODO: 总结本章内容

}

